\documentclass[12pt]{article}
\usepackage[utf8]{inputenc}
\usepackage[T1]{fontenc}
\usepackage{hyperref}
\title{Weh Language Resources for K-12 Curriculum Development}
\author{Compiled by Claude Sonnet 5 (claude-sonnet-5) for Andr\'as Kornai}
\date{}
\begin{document}
\maketitle
\section{Introduction}
Weh (ISO 639-3 weh; autonym K\'\i w\'\i) is a Grassfields Bantu language of the Menchum Division in Cameroon's North West Region. Speaker counts vary by source, from about 7,000 to 10,400.
\section{Primary Educational Resources}
None found in search results.
\section{Linguistic and Reference Materials}
\begin{itemize}
\item Weh orthography guide (SIL Cameroon, 2009): \url{https://www.silcam.org/resources/archives/47887}
\item Weh phonology sketch (SIL Cameroon, 2009): \url{https://www.silcam.org/resources/archives/47888}
\item Glottolog reference for Davison (2004), an earlier orthography guide: \url{https://glottolog.org/resource/reference/id/84218}
\item Omniglot page on Weh: \url{https://www.omniglot.com/writing/weh.htm}
\end{itemize}
\section{Religious and Community Resources}
Joshua Project reports a New Testament (2022): \url{https://joshuaproject.net/languages/weh}. I did not find the text itself.
\section{Digital and Audio Resources}
None found in search results.
\section{Curriculum Development Context}
An orthography proposal and a phonology sketch exist, which is a basis for materials. The orthography needs further testing, and primers and readers remain to be developed.
\end{document}